\section{Post-training y compute en tiempo de inferencia: cómo el comportamiento continúa cambiando después del preentrenamiento}

El \texttt{pretraining} permite al modelo adquirir una distribución condicional amplia, pero no convierte automáticamente las salidas en un asistente controlable. El \texttt{post-training} puede modificar los parámetros o añadir durante la inferencia prompts, búsqueda, recuperación, herramientas o verificadores. Al evaluar un método, pregúntese primero: ¿actualiza pesos, modifica el decodificación, añade estado externo o redefine las recompensas? De lo contrario, CoT, RLHF, RAG y agent corren el riesgo de mezclarse en una capa de "mejora" difusa.

\subsection{Mapa del post-training: ¿qué más hay que ajustar después del entrenamiento}

El curso coloca \texttt{post-training} y \texttt{inference-time compute} en la misma página: \texttt{prompt engineering}, ajuste de instrucciones (\texttt{instruction tuning}), optimización de preferencias, recuperación, agentes/herramientas pueden todos cambiar el comportamiento desplegado. Si la distribución del modelo base es $p_0(y\mid x)$, el post-training produce $p_{\theta}(y\mid x)$, y la estrategia en tiempo de inferencia puede escribirse como
\[
y^{\star}=\arg\max_{y\in\mathcal{S}(x;B)}
\left[s_{\theta}(y\mid x)+\lambda r(y,x)\right],
\]
donde $\mathcal{S}(x;B)$ es el conjunto de candidatos generados dentro del presupuesto $B$, $s_{\theta}$ es la puntuación del modelo y $r$ es el verificador/recompensa. Aumentar el compute en tiempo de inferencia expande el conjunto de candidatos o la profundidad de la verificación, pero no repara conocimientos completamente ausentes ni permisos de herramientas faltantes.

\lecturefigure{slide-088.jpg}{Mapa de métodos de post-training e inference-time compute}{CS25 V6 oficial, p.~88}

\teachervoice{00:36:33--00:38:05，el docente considera prompting, fine-tuning, aprendizaje de preferencias, recuperación y uso de agentes/herramientas como diferentes puntos de intervención. Nota de clase: el comportamiento desplegado surge de la interacción entre el modelo base y el sistema de inferencia.}

\begin{importantbox}{Cuatro significados de "mejorar un modelo"}
\begin{enumerate}
  \item \textbf{Actualizar pesos}: SFT, DPO, RL, etc., cambian directamente los parámetros;
  \item \textbf{Modificar decodificación}: autoconsistencia, búsqueda y reranking alteran la selección de candidatos;
  \item \textbf{Añadir contexto/estado}: RAG, memoria y salidas de herramientas añaden información visible;
  \item \textbf{Cambiar el entorno}: el ciclo del agente permite al modelo actuar y obtener nuevas observaciones.
\end{enumerate}
Los costos, la capacidad de reversión y los patrones de fallo de las cuatro son completamente distintos.
\end{importantbox}

\subsection{De CoT a descomposición ejecutable}

Chain-of-Thought (CoT) exige que el modelo produzca pasos intermedios, dividiendo una mapeo único en razonamiento secuencial. Puede mejorar el rendimiento en tareas complejas o generar explicaciones que suenen lógicas pero sean infieles. Tree-of-Thoughts amplía esto a una búsqueda de múltiples ramas; Program-of-Thought delega parte del cálculo a un intérprete; Socratic questioning y graph de cómputo definen explícitamente las dependencias entre subproblemas.

\lecturefigure{slide-089.jpg}{Chain-of-Thought: permitir al modelo desplegar pasos intermedios antes de la respuesta}{CS25 V6 oficial, p.~89}

\lecturefigure{slide-090.jpg}{Comparativa de ejemplos entre prompting estándar y CoT prompting}{CS25 V6 oficial, p.~90}

\lecturefigure{slide-091.jpg}{Tree-of-Thoughts: expansión desde una cadena de razonamiento única a búsqueda ramificada}{CS25 V6 oficial, p.~91}

\lecturefigure{slide-092.jpg}{Program-of-Thought: delegar cómputos ejecutables al intérprete de código}{CS25 V6 oficial, p.~92}

\begin{knowledgebox}{Lectura de la figura: los cuatro métodos alteran interfaces distintas}
CoT añade tokens secuenciales; ToT añade ramas de candidatos y un controlador de búsqueda; Program-of-Thought añade una herramienta ejecutable; Socratic/graph de cómputo añaden una estructura de dependencia explícita. Todos pueden llamarse "razonamiento", pero difieren en compute, verificabilidad y rutas de propagación de errores.
\end{knowledgebox}

Socratic decomposition puede escribirse como un grafo acíclico dirigido $G=(V,E)$, donde cada nodo $v_i$ es un subproblema y la arista $(v_j,v_i)$ indica que $i$ depende de $j$. Solo se resuelve $v_i$ cuando todos sus predecesores están completos. Esta estructura facilita la verificación local, pero introduce error de descomposición: si se omiten nodos clave al inicio, cada paso posterior puede ser localmente correcto y globalmente fallido.

\lecturefigure{slide-093.jpg}{Socratic questioning: preguntas recursivas que dividen problemas complejos en subproblemas}{CS25 V6 oficial, p.~93}

\lecturefigure{slide-094.jpg}{Graphs de cómputo: organizar descomposición y resultados intermedios mediante dependencias}{CS25 V6 oficial, p.~94}

\teachervoice{00:38:05--00:40:36，el docente divide las extensiones de CoT en branching, ejecución y descomposición estructurada. Recordatorio del material: más tokens de razonamiento son presupuesto de búsqueda, no equivalen a que esos tokens reflejen fielmente el proceso causal interno del modelo.}

\begin{lstlisting}[language=Python,caption={Pseudocódigo de búsqueda de razonamiento estructurado con verificador}]
frontier = [initial_state(problem)]
for depth in range(max_depth):
    candidates = []
    for state in frontier:
        for thought in model.propose(state, branches=k):
            next_state = state.extend(thought)
            score = verifier(next_state)
            candidates.append((score, next_state))
    frontier = top_b(candidates, beam_width)
return select_final(frontier)
\end{lstlisting}

\begin{warningbox}{Tres niveles de evidencia para la reasoning trace}
Una trazabilidad puede ser una \textbf{explicación legible}, un \textbf{scratchpad que ayuda a obtener la respuesta} o una \textbf{evidencia causal fiel del cálculo interno}. Los dos primeros son comunes; el tercero requiere intervención, contrafactuales o evidencia mecánica. No se debe afirmar que "el modelo pensó exactamente así" solo porque los pasos textuales sean correctos.
\end{warningbox}

\subsection{RLHF, DPO, RLAIF y GRPO: fuentes de retroalimentación y rutas de optimización distintas}

Reinforcement Learning from Human Feedback (RLHF) típicamente recopila preferencias primero, entrena un modelo de recompensa y luego usa optimización de política para cambiar la estrategia de generación. Direct Preference Optimization (DPO) salta el modelo de recompensa explícito y el RL online, optimizando directamente los pares de preferencias:
\[
\mathcal{L}_{\mathrm{DPO}}
=-\log\sigma\!\left(\beta\left[
\log\frac{\pi_{\theta}(y^+\mid x)}{\pi_{\mathrm{ref}}(y^+\mid x)}
-\log\frac{\pi_{\theta}(y^-\mid x)}{\pi_{\mathrm{ref}}(y^-\mid x)}
\right]\right).
\]
$y^+$ y $y^-$ son respuestas preferidas/rechazadas, $\pi_{\mathrm{ref}}$ es la política de referencia y $\beta$ controla el grado de desviación.

\lecturefigure{slide-096.jpg}{RLHF: preferencias humanas, modelo de recompensa y optimización de política}{CS25 V6 oficial, p.~96}

\lecturefigure{slide-097.jpg}{DPO: actualizar la política directamente desde pares de preferencias sin entrenar explícitamente un modelo de recompensa}{CS25 V6 oficial, p.~97}

RLAIF usa IA para generar o juzgar preferencias, reduciendo costos manuales, pero puede amplificar los sesgos del modelo maestro. GRPO muestrea un conjunto de respuestas bajo el mismo prompt y construye una ventaja relativa dentro del grupo. Si la recompensa es $r_i$, con media y desviación estándar $\mu_r,\sigma_r$ del grupo, se puede escribir como
\[
\hat{A}_i=\frac{r_i-\mu_r}{\sigma_r+\epsilon}.
\]
Evita el costo de un modelo value independiente, pero aún depende de la fiabilidad del verificador/recompensa y puede generar altos costos computacionales mediante muestreo por grupos.

\lecturefigure{slide-098.jpg}{RLAIF: retroalimentación de IA sustituye parcialmente la anotación manual de preferencias}{CS25 V6 oficial, p.~98}

\lecturefigure{slide-099.jpg}{GRPO: recompensa relativa intra-grupo y optimización de política de razonamiento}{CS25 V6 oficial, p.~99}

\begin{knowledgebox}{Términos clave: ¿en qué difieren realmente los métodos de optimización de preferencias?}
\begin{tabularx}{\textwidth}{p{0.17\textwidth}p{0.24\textwidth}X}
\toprule
Método & Fuente de retroalimentación/modelo intermedio & Mecánica central y riesgos principales \\
\midrule
RLHF & Preferencias humanas + modelo de recompensa & Optimización de política online/casi-online; hacking de recompensas y alta complejidad de entrenamiento \\
DPO & Pares de preferencias humanos o IA & Objetivo tipo clasificación directa; más simple pero aún sujeto a la calidad de los datos de preferencia \\
RLAIF & Crítico de IA o retroalimentación constitucional & Reduce costos manuales; el sesgo del maestro puede ser amplificado sistemáticamente \\
GRPO & Recompensa verificable + muestras intra-grupo & Ventaja relativa, sin modelo value independiente; costo de muestreo y sesgo del verificador son evidentes \\
\bottomrule
\end{tabularx}
\end{knowledgebox}

\subsection{Supervisión del proceso: recompensar pasos o solo el resultado}

La supervisión de resultados (Outcome supervision) mira solo la respuesta final; el modelo puede obtener puntos accidentalmente mediante atajos. La supervisión del proceso (process supervision) proporciona etiquetas o recompensas para los pasos intermedios. Si la trazabilidad de razonamiento es $z_{1:m}$, se puede escribir como
\[
R_{\mathrm{process}}(z)=\sum_{i=1}^{m}w_i r_i(z_i,z_{<i},x).
\]
$r_i$ evalúa el paso $i$, y $w_i$ es el peso. La dificultad radica en quién define los "pasos correctos", cómo manejar múltiples rutas equivalentes y si el modelo aprende a complacer al verificador en lugar de resolver el problema.

\lecturefigure{slide-100.jpg}{Supervisión del proceso: pasar de supervisar la respuesta final a supervisar pasos intermedios de razonamiento}{CS25 V6 oficial, p.~100}

\teachervoice{00:40:36--00:43:12，el docente distingue RLHF, DPO, RLAIF, GRPO y supervisión del proceso: difieren en la fuente de retroalimentación, si construyen explícitamente un modelo de recompensa y si recompensan el resultado o los pasos.}

\begin{warningbox}{Las tareas verificables no implican que el verificador sea imparcial}
Las respuestas matemáticas, las pruebas de código y las verificaciones de formato se pueden calificar automáticamente fácilmente, pero la recompensa a menudo cubre solo partes medibles. El modelo puede generar respuestas sobreajustadas a las pruebas, aprovechar tolerancias numéricas o producir "pasos calificables" redundantes. Por tanto, el diseño del verificador es en sí mismo un objetivo de entrenamiento, no una herramienta de medición neutral.
\end{warningbox}

Desde la infraestructura, DPO requiere principalmente pares de preferencias y retropropagación estándar; RLHF/GRPO requieren generación de rollout, cálculo de recompensa, logprob de política/modelo de referencia y pueden intercambiar secuencias de longitud variable masiva entre múltiples máquinas. Los tokens de razonamiento también son costo de entrenamiento: las diferencias en longitud de razonamiento dentro del mismo batch causan padding y carga desigual. Al elegir un algoritmo, se debe comparar simultáneamente la construcción de datos offline, muestreo online, latencia de recompensa, copias múltiples de modelos en memoria y reproducibilidad. Que un algoritmo tenga mayor eficiencia de tokens en un artículo no significa que sea mejor en tiempo real (wall-clock) o uso del clúster.

\subsection{Resumen de la sección}

El post-training puede modificar parámetros, el compute en tiempo de inferencia puede expandir la búsqueda, la recuperación/herramientas pueden añadir información y la supervisión del proceso puede cambiar la granularidad de la recompensa. CoT, ToT, Program-of-Thought, RLHF, DPO, RLAIF y GRPO no son una misma clase de "truco de razonamiento"; operan sobre espacios de candidatos, ejecución externa, datos de preferencia y optimizadores respectivamente.
